% NB15b -- Condensador de superficie (coraza y tubos, 2 pasos) y su red de capacitancias concentradas
% (vapor -> pelicula -> pared -> agua de enfriamiento) que define las EDO en el tiempo.
% Reemplaza "CORAZA Y TUBOS.png" de la entrega PlantaTermoelectrica (Sanchez, Aparicio, Roa).
\documentclass[tikz,border=4mm]{standalone}
\usepackage{amsmath,amssymb}
\usetikzlibrary{arrows.meta,calc,patterns}
\begin{document}
% Estilo comun de los diagramas TikZ del modulo VCR (NB10b) -- paleta del proyecto
\definecolor{manivela}{HTML}{2A78D6}
\definecolor{inferior}{HTML}{4A3AA7}
\definecolor{control}{HTML}{1BAF7A}
\definecolor{superior}{HTML}{EB6834}
\definecolor{tinta}{HTML}{3A3A38}
\definecolor{metal}{HTML}{C3C2B7}
\definecolor{perdida}{HTML}{E34948}
\definecolor{fuerza}{HTML}{C0392B}
\tikzset{
  barra/.style={line width=3.2pt, line cap=round},
  pin/.style={circle, fill=white, draw=tinta, line width=0.9pt, inner sep=1.6pt},
  fijo/.style={circle, fill=tinta, inner sep=2.6pt},
  vec/.style={-{Stealth[length=2.6mm]}, line width=0.9pt},
  F/.style={-{Stealth[length=3mm]}, line width=1.2pt, fuerza},
  cota/.style={{Stealth[length=1.8mm]}-{Stealth[length=1.8mm]}, line width=0.5pt, tinta},
  lbl/.style={font=\small, text=tinta},
}
\newcommand{\tierra}[1]{\fill[pattern=north east lines, pattern color=tinta] ($(#1)+(-5,-3)$) rectangle ($(#1)+(5,-6)$);
  \draw[tinta, line width=0.8pt] ($(#1)+(-5,-3)$) -- ($(#1)+(5,-3)$);}

\definecolor{frio}{HTML}{2A78D6}
\definecolor{vapor}{HTML}{4A3AA7}
\definecolor{caliente}{HTML}{E34948}
\begin{tikzpicture}[x=1cm,y=1cm,
  flujo/.style={-{Stealth[length=2.6mm]}, line width=1.6pt},
  nodo/.style={circle, draw=tinta, line width=1pt, minimum size=13mm, align=center, font=\footnotesize},
  res/.style={draw=tinta, line width=0.8pt, fill=white, minimum width=9mm, minimum height=3.5mm, font=\scriptsize}]
  \node[lbl, font=\small\bfseries, anchor=west] at (-0.6,5.4) {(a) Condensador de superficie, 2 pasos por tubos};
  % coraza
  \draw[tinta, line width=1.2pt, fill=vapor!6, rounded corners=14pt] (0,0) rectangle (9,4.2);
  % placas tubulares y cabezales
  \draw[tinta, line width=1pt] (1.0,0) -- (1.0,4.2); \draw[tinta, line width=1pt] (8.0,0) -- (8.0,4.2);
  \fill[frio!18] (0.12,0.12) rectangle (0.98,4.08); \fill[frio!18] (8.02,0.12) rectangle (8.88,4.08);
  \draw[tinta, line width=1pt] (0.1,2.1) -- (1.0,2.1); % particion de pasos
  % tubos: paso 1 abajo, paso 2 arriba
  \foreach \y in {0.6,1.0,1.4,1.8} \draw[frio!80!tinta, line width=1.4pt] (1.0,\y) -- (8.0,\y);
  \foreach \y in {2.4,2.8,3.2,3.6} \draw[frio!40!caliente, line width=1.4pt] (1.0,\y) -- (8.0,\y);
  \foreach \x in {2.8,5.2} \draw[tinta, line width=0.9pt] (\x,4.2) -- (\x,1.1);
  \foreach \x in {4.0,6.4} \draw[tinta, line width=0.9pt] (\x,0.0) -- (\x,3.1);
  \node[lbl, font=\scriptsize, fill=white, inner sep=1pt] at (2.8,4.42) {deflector};
  % entrada de vapor y pozo caliente
  \draw[flujo, vapor] (4.6,5.0) -- (4.6,4.2); \node[lbl, font=\footnotesize, vapor, left] at (4.5,4.85) {vapor de escape $\dot m_s,\;h_4$ (4)};
  \fill[frio!45] (3.3,-0.6) rectangle (5.9,0.0); \draw[tinta, line width=1pt] (3.3,0) -- (3.3,-0.6) -- (5.9,-0.6) -- (5.9,0);
  \node[lbl, font=\scriptsize] at (4.6,-0.32) {pozo caliente};
  \draw[flujo, frio] (4.6,-0.6) -- (4.6,-1.3); \node[lbl, font=\footnotesize, right] at (4.7,-1.05) {condensado $\dot m_{cd}$ a la bomba (1)};
  % venteo
  \draw[flujo, tinta!60] (7.4,4.2) -- (7.4,4.9); \node[lbl, font=\scriptsize, right] at (7.5,4.75) {venteo (no condensables)};
  % agua de enfriamiento
  \draw[flujo, frio] (-1.2,1.2) -- (0.1,1.2); \node[lbl, font=\scriptsize, align=right, anchor=east] at (-1.25,1.2) {$T_{cw,in}$};
  \draw[flujo, frio!40!caliente] (0.1,3.0) -- (-1.2,3.0); \node[lbl, font=\scriptsize, align=right, anchor=east] at (-1.25,3.0) {$T_{cw,out}$};
  \draw[-{Stealth[length=2mm]}, tinta!70, line width=0.8pt] (8.45,1.4) arc (-90:90:0.35 and 0.95);
  \node[lbl, font=\scriptsize, align=center] at (4.5,-2.0) {$N_t=3000$ tubos, $L=10$ m, $D_o=25$ mm (geometr\'ia de la entrega)};
  % ================= red de capacitancias =================
  \begin{scope}[shift={(11.3,0)}]
    \node[lbl, font=\small\bfseries, anchor=west] at (-0.5,5.4) {(b) Red t\'ermica concentrada $\Rightarrow$ 3 balances};
    \node[nodo, fill=vapor!15] (V) at (1,3.6) {vapor\\$p_c,\,M_v$};
    \node[nodo, fill=metal!40] (Wl) at (1,1.0) {pared\\$T_w$};
    \node[nodo, fill=frio!20] (A) at (4.6,1.0) {agua\\$\bar T_{cw}$};
    \draw[tinta, line width=1pt] (V) -- node[res, pos=0.5] {$h_oA_o$} (Wl);
    \draw[tinta, line width=1pt] (Wl) -- node[res, pos=0.5] {$G_{eq}$} (A);
    \draw[-{Stealth[length=2mm]}, vapor, line width=1pt] (-0.4,4.6) -- (V) node[pos=0, left, font=\scriptsize] {$\dot m_s h_4$};
    \draw[-{Stealth[length=2mm]}, frio, line width=1pt] (V.east) -- ++(1.6,0) node[right, font=\scriptsize, align=left] {$\dot m_{cd}=\dfrac{h_oA_o(T_{sat}-T_w)}{h_{fg}}$};
    \draw[-{Stealth[length=2mm]}, frio, line width=1pt] (4.6,-0.4) -- (A) node[pos=0, below, font=\scriptsize] {$\dot m_{cw}c_pT_{cw,in}$};
    \node[draw=tinta, rounded corners=2pt, fill=white, font=\footnotesize, align=left, anchor=north west] at (-0.6,-0.9)
      {$V_c\,\dfrac{d\rho_g}{dp}\Big|_{sat}\dfrac{dp_c}{dt}=x_4\dot m_s-\dot m_{cd}$\quad ($\tau_p\sim$ ms $\Rightarrow$ cuasi-est\'atico: $x_4\dot m_s=\dot m_{cd}$)\\[3pt]
       $C_w\dfrac{dT_w}{dt}=h_oA_o\,(T_{sat}-T_w)-G_{eq}\,(T_w-\bar T_{cw})$\\[3pt]
       $C_{cw}\dfrac{d\bar T_{cw}}{dt}=G_{eq}\,(T_w-\bar T_{cw})-\dot m_{cw}c_p\,(T_{cw,out}-T_{cw,in})$\\[3pt]
       $G_{eq}$: $h_iA_i$ corregido para reproducir $\varepsilon$-NTU;\quad $\tau_w=\dfrac{C_w}{G_{eq}+h_oA_o}$,\quad $\tau_{cw}=\dfrac{C_{cw}}{G_{eq}+2\dot m_{cw}c_p}$ (calculadas)};
  \end{scope}
\end{tikzpicture}
\end{document}
